% !TeX root = ../../main.tex

\section{\bilingual{Random-Walk Lemmas}{随机游走引理}}
\label{section:random-walk-lemmas}

\begin{lemma}[\bilingual{Meander estimate for the transition event}{基于 Brownian 蜿蜒的过渡事件估计}]
    \label{lemma:meander-transition}
    \bilingual{
        Let $(\xi_j)_{j\ge1}$ be i.i.d. random variables with mean $0$ and variance $\sigma^2\in(0,\infty)$. Fix $c\ge0$. For $\Delta>0$, set $n_\Delta:=\lfloor\Delta^2\rfloor$ and $h_j:=cj/\Delta^2$. For $x\in[0,\Delta]$, let $S_0=x$, $S_j:=x+\sum_{i=1}^j\xi_i$, and define
    }{
        设 $(\xi_j)_{j\ge1}$ 为均值 $0$、方差 $\sigma^2\in(0,\infty)$ 的独立同分布随机变量. 固定 $c\ge0$. 对 $\Delta>0$, 令 $n_\Delta:=\lfloor\Delta^2\rfloor$ 且 $h_j:=cj/\Delta^2$. 对 $x\in[0,\Delta]$, 令 $S_0=x$, $S_j:=x+\sum_{i=1}^j\xi_i$, 并定义
    }
    \begin{equation*}
        \mathcal A_\Delta(x):=
        \left\{
            h_j\le S_j\le\Delta+h_j\ (1\le j\le n_\Delta),\quad
            S_{n_\Delta}\in\left[h_{n_\Delta}+\frac25\Delta,
                h_{n_\Delta}+\frac35\Delta\right]
        \right\}.
    \end{equation*}
    \bilingual{Then there are $C>0$ and $\Delta_0>0$ such that}{则存在 $C>0$ 与 $\Delta_0>0$, 使得}
    \begin{equation*}
        \inf_{x\in[0,\Delta]}\mathbb P_x\bigl(\mathcal A_\Delta(x)\bigr)
        \ge C\Delta^{-2},
        \qquad \Delta\ge\Delta_0,
    \end{equation*}
    \bilingual{
        Here $\mathbb P_x$ denotes the law with $S_0=x$.
    }{
        其中 $\mathbb P_x$ 表示 $S_0=x$ 时的概率律.
    }
\end{lemma}

\begin{proof}[\bilingual{Proof of \Cref{lemma:meander-transition}}{\Cref{lemma:meander-transition} 的证明}]
    \bilingual{
        For $\theta\in\{-1,1\}$, let $Z_0^\theta=0$ and let the increments of $(Z_j^\theta)$ have the law of $\theta\xi_1$. Put $C_N^\theta:=\{Z_j^\theta\ge0,\ 1\le j\le N\}$ and define
    }{
        对 $\theta\in\{-1,1\}$, 令 $Z_0^\theta=0$, 且 $(Z_j^\theta)$ 的增量服从 $\theta\xi_1$ 的分布. 记 $C_N^\theta:=\{Z_j^\theta\ge0,\ 1\le j\le N\}$, 并定义
    }
    \begin{equation*}
        D_{N,\Delta}^\theta:=\left\{
            0\le Z_j^\theta\le\frac56\Delta\ (1\le j\le N),\quad
            Z_N^\theta\in\left[\frac5{12}\Delta,\frac{11}{24}\Delta\right]
        \right\}.
    \end{equation*}
    \bilingual{
        In the Skorokhod space $D([0,1])$ of càdlàg paths, the conditioned invariance principle of Bolthausen~\mycite{Bolthausen_1976}, Theorem 3.2, together with the meander identification in Iglehart~\mycite{Iglehart_1974}, Section 3, and Caravenna and Chaumont~\mycite{Caravenna_Chaumont_2008}, equation (3.3), gives
    }{
        在配备 Skorokhod 拓扑的 càdlàg 路径空间 $D([0,1])$ 中, 由 Bolthausen~\mycite{Bolthausen_1976} 的定理 3.2、Iglehart~\mycite{Iglehart_1974} 的第 3 节及 Caravenna 与 Chaumont~\mycite{Caravenna_Chaumont_2008} 的公式 (3.3) 中对 Brownian 蜿蜒的识别, 得到条件不变原理
    }
    \begin{equation*}
        \mathcal L\left(\left.\left(\frac{Z_{\lfloor Nt\rfloor}^\theta}{\sigma\sqrt N}\right)_{0\le t\le1}\,\right|C_N^\theta\right)
        \Longrightarrow \mathsf m,
    \end{equation*}
    \bilingual{
        where $\mathsf m$ is the standard Brownian meander. Fix $\rho:=\frac1{100}$ and let
    }{
        其中 $\mathsf m$ 为标准 Brownian 蜿蜒. 令 $\rho:=\frac1{100}$, 并定义
    }
    \begin{equation*}
        \mathcal U:=\left\{f\in D([0,1]):
            \sigma\sup_{0\le t\le1}f(t)<\frac56-\rho,\quad
            \sigma f(1)\in\left(\frac5{12}+\rho,\frac{11}{24}-\rho\right)
        \right\}.
    \end{equation*}
    \bilingual{
        The set $\mathcal U$ is open in the Skorokhod topology, and the support of the Brownian meander contains the path $f_0(t)=\frac{7t}{16\sigma}\in\mathcal U$. Thus $q_*:=\mathsf m(\mathcal U)>0$. For each fixed integer $m\ge0$, put $N=n_\Delta-m$; Portmanteau and $\Delta/\sqrt N\to1$ yield
    }{
        $\mathcal U$ 在 Skorokhod 拓扑下是开集, Brownian 蜿蜒的支撑包含路径 $f_0(t)=\frac{7t}{16\sigma}\in\mathcal U$. 故 $q_*:=\mathsf m(\mathcal U)>0$. 对每个固定整数 $m\ge0$, 令 $N=n_\Delta-m$; 由 Portmanteau 定理及 $\Delta/\sqrt N\to1$, 得
    }
    \begin{equation*}
        \liminf_{\Delta\to\infty}
        \mathbb P\left(D_{n_\Delta-m,\Delta}^\theta\middle|C_{n_\Delta-m}^\theta\right)
        \ge q_*.
    \end{equation*}
    \bilingual{
        To bound the survival probability, write $F_N^\theta:=\mathbb P(C_N^\theta)$ and $p_N^\theta:=\mathbb P(Z_N^\theta\ge0)$. On $\{Z_N^\theta\ge0\}$, the cyclic rotation cover of Bauer, Godrèche and Luck~\mycite{Bauer_Godreche_Luck_1999}, Appendix A, following (A.3), gives a rotation whose partial sums are nonnegative. The $N$ rotated increment vectors have the same law by the i.i.d. assumption. The union bound therefore gives
    }{
        为估计存活概率, 令 $F_N^\theta:=\mathbb P(C_N^\theta)$ 与 $p_N^\theta:=\mathbb P(Z_N^\theta\ge0)$. 在事件 $\{Z_N^\theta\ge0\}$ 上, Bauer、Godrèche 与 Luck~\mycite{Bauer_Godreche_Luck_1999} 附录 A 公式 (A.3) 后的循环移位覆盖保证存在一个所有部分和均非负的移位. 由独立同分布性, $N$ 个移位后的增量向量具有相同的分布. 对这些事件应用并集上界, 得
    }
    \begin{equation*}
        p_N^\theta\le NF_N^\theta,
        \qquad N\ge1.
    \end{equation*}
    \bilingual{
        By the central limit theorem, $p_N^\theta\to\frac12$. Thus $p_N^\theta\ge\frac14$ for all sufficiently large $N$, and the preceding inequality gives
    }{
        由中心极限定理, $p_N^\theta\to\frac12$. 因此充分大的 $N$ 满足 $p_N^\theta\ge\frac14$, 上述不等式给出
    }
    \begin{equation*}
        F_N^\theta\ge\frac{1}{4N}.
    \end{equation*}
    \bilingual{
        For each fixed integer $m\ge0$, $N=n_\Delta-m\to\infty$ and $N\le\Delta^2$. Consequently $F_N^\theta\ge\frac{1}{4\Delta^2}$ for all sufficiently large $\Delta$. Since $D_{N,\Delta}^\theta\subseteq C_N^\theta$, the conditional estimate above and $\mathbb P(D_{N,\Delta}^\theta)=\mathbb P(D_{N,\Delta}^\theta\mid C_N^\theta)F_N^\theta$ give $c_m:=q_*/8>0$ such that
    }{
        对每个固定整数 $m\ge0$, 有 $N=n_\Delta-m\to\infty$ 且 $N\le\Delta^2$. 因此充分大的 $\Delta$ 满足 $F_N^\theta\ge\frac{1}{4\Delta^2}$. 由 $D_{N,\Delta}^\theta\subseteq C_N^\theta$, 结合上述条件概率估计及 $\mathbb P(D_{N,\Delta}^\theta)=\mathbb P(D_{N,\Delta}^\theta\mid C_N^\theta)F_N^\theta$, 取 $c_m:=q_*/8>0$, 得
    }
    \begin{equation*}
        \min_{\theta\in\{-1,1\}}
        \mathbb P\left(D_{n_\Delta-m,\Delta}^\theta\right)
        \ge c_m\Delta^{-2}
    \end{equation*}
    \bilingual{for all sufficiently large $\Delta$.}{对充分大的 $\Delta$ 成立.}

    \bilingual{
        For starting points in the middle interval, Brownian path support and compactness give
    }{
        对中间区间内的起点, Brownian 路径的支撑性质与紧致性给出
    }
    \begin{equation*}
        p_0:=\inf_{z\in[\frac1{12},\frac{11}{12}]}
        \mathbb P\left(
            \frac1{24}<z+\sigma B_t<\frac{23}{24}\ (0\le t\le1),\quad
            z+\sigma B_1\in\left(\frac5{12},\frac7{12}\right)
        \right)>0,
    \end{equation*}
    \bilingual{
        where $B$ is standard Brownian motion. Since $n_\Delta/\Delta^2\to1$ and $c/\Delta\to0$, Donsker's invariance principle and compactness of $[\frac1{12},\frac{11}{12}]$ imply
    }{
        其中 $B$ 为标准 Brownian 运动. 由于 $n_\Delta/\Delta^2\to1$ 且 $c/\Delta\to0$, 由 Donsker 不变原理及区间 $[\frac1{12},\frac{11}{12}]$ 的紧致性, 得
    }
    \begin{equation*}
        \inf_{x\in[\frac1{12}\Delta,\frac{11}{12}\Delta]}
        \mathbb P_x\left(
            c\le S_j\le\Delta\ (1\le j\le n_\Delta),\quad
            S_{n_\Delta}\in\left[\frac5{12}\Delta,\frac7{12}\Delta\right]
        \right)\ge\frac{p_0}{2}
    \end{equation*}
    \bilingual{for sufficiently large $\Delta$.}{对充分大的 $\Delta$ 成立.}

    \bilingual{
        The two edge intervals use a fixed initial segment followed by a meander segment. Since the increment law is nondegenerate and centered, choose $0<a_+<b_+$ and $0<a_-<b_-$ such that
    }{
        两侧区间均先取固定长度的初始段, 再接 Brownian 蜿蜒段. 由于增量分布非退化且均值为零, 可取 $0<a_+<b_+$ 与 $0<a_-<b_-$ 使得
    }
    \begin{equation*}
        p_+:=\mathbb P(\xi_1\in[a_+,b_+])>0,
        \qquad
        p_-:=\mathbb P(-\xi_1\in[a_-,b_-])>0.
    \end{equation*}
    \bilingual{
        Fix integers $m_+,m_-$ with $m_+a_+>c+1$ and $m_-a_->c+1$. For $x\le\frac1{12}\Delta$, impose $\xi_i\in[a_+,b_+]$ for $1\le i\le m_+$ and use $D_{n_\Delta-m_+,\Delta}^{1}$ on the independent suffix. For sufficiently large $\Delta$, the prefix satisfies $h_j<S_j\le\Delta$ for $1\le j\le m_+$, and the suffix satisfies $c<S_{m_+}+Z_j^1\le\Delta$ for $1\le j\le n_\Delta-m_+$. At the endpoint,
    }{
        固定整数 $m_+,m_-$ 使得 $m_+a_+>c+1$ 且 $m_-a_->c+1$. 对 $x\le\frac1{12}\Delta$, 在 $1\le i\le m_+$ 时要求 $\xi_i\in[a_+,b_+]$, 并在独立后缀上取事件 $D_{n_\Delta-m_+,\Delta}^{1}$. 对充分大的 $\Delta$, 前缀满足 $h_j<S_j\le\Delta$, $1\le j\le m_+$; 后缀满足 $c<S_{m_+}+Z_j^1\le\Delta$, $1\le j\le n_\Delta-m_+$. 在终点,
    }
    \begin{equation*}
        h_{n_\Delta}+\frac25\Delta
        \le\frac5{12}\Delta\le S_{n_\Delta}
        \le\frac{13}{24}\Delta+m_+b_+
        \le h_{n_\Delta}+\frac35\Delta.
    \end{equation*}
    \bilingual{
        Indeed, the upper bound for the suffix is $x+m_+b_++\frac56\Delta\le\frac{11}{12}\Delta+m_+b_+\le\Delta$, while the endpoint lies in $[\frac5{12}\Delta,\frac{13}{24}\Delta+m_+b_+]$. Independence and the meander estimate therefore give
    }{
        事实上, 后缀的上界为 $x+m_+b_++\frac56\Delta\le\frac{11}{12}\Delta+m_+b_+\le\Delta$, 终点位于 $[\frac5{12}\Delta,\frac{13}{24}\Delta+m_+b_+]$. 由独立性及 Brownian 蜿蜒估计, 得
    }
    \begin{equation*}
        \inf_{0\le x\le\frac1{12}\Delta}\mathbb P_x(\mathcal A_\Delta(x))
        \ge p_+^{m_+}c_{m_+}\Delta^{-2}.
    \end{equation*}

    \bilingual{
        For $x\ge\frac{11}{12}\Delta$, set $R_j:=\Delta-S_j$ and $y:=\Delta-x$. Require $-\xi_i\in[a_-,b_-]$ for $1\le i\le m_-$, and use $D_{n_\Delta-m_-,\Delta}^{-1}$ on the independent suffix. For sufficiently large $\Delta$, the prefix satisfies $0\le R_j\le\Delta-c$ for $1\le j\le m_-$, and the suffix satisfies $0\le R_{m_-}+Z_j^{-1}\le\Delta-c$ for $1\le j\le n_\Delta-m_-$. Consequently $c\le S_j\le\Delta$ for $1\le j\le n_\Delta$. At the endpoint,
    }{
        对 $x\ge\frac{11}{12}\Delta$, 令 $R_j:=\Delta-S_j$ 且 $y:=\Delta-x$. 在 $1\le i\le m_-$ 时要求 $-\xi_i\in[a_-,b_-]$, 并在独立后缀上取事件 $D_{n_\Delta-m_-,\Delta}^{-1}$. 对充分大的 $\Delta$, 前缀满足 $0\le R_j\le\Delta-c$, $1\le j\le m_-$; 后缀满足 $0\le R_{m_-}+Z_j^{-1}\le\Delta-c$, $1\le j\le n_\Delta-m_-$. 因此 $c\le S_j\le\Delta$, $1\le j\le n_\Delta$. 在终点,
    }
    \begin{equation*}
        h_{n_\Delta}+\frac25\Delta
        \le\frac{11}{24}\Delta-m_-b_-
        \le S_{n_\Delta}\le\frac7{12}\Delta
        \le h_{n_\Delta}+\frac35\Delta.
    \end{equation*}
    \bilingual{
        Here $R_{n_\Delta}\le\frac{13}{24}\Delta+m_-b_-$ gives $S_{n_\Delta}\ge\frac{11}{24}\Delta-m_-b_-$, and $R_{n_\Delta}\ge\frac5{12}\Delta$ gives $S_{n_\Delta}\le\frac7{12}\Delta$. Independence yields
    }{
        其中 $R_{n_\Delta}\le\frac{13}{24}\Delta+m_-b_-$ 给出 $S_{n_\Delta}\ge\frac{11}{24}\Delta-m_-b_-$, 而 $R_{n_\Delta}\ge\frac5{12}\Delta$ 给出 $S_{n_\Delta}\le\frac7{12}\Delta$. 由独立性得
    }
    \begin{equation*}
        \inf_{\frac{11}{12}\Delta\le x\le\Delta}\mathbb P_x(\mathcal A_\Delta(x))
        \ge p_-^{m_-}c_{m_-}\Delta^{-2}.
    \end{equation*}

    \bilingual{
        For the middle interval, $c\le S_j\le\Delta$ implies $h_j\le S_j\le\Delta+h_j$, because $h_j\le c$ for $j\le n_\Delta$. Since $h_{n_\Delta}\le c$ and $\frac25<\frac5{12}<\frac7{12}<\frac35$, its endpoint interval lies in the one defining $\mathcal A_\Delta(x)$ for sufficiently large $\Delta$. The displayed endpoint inequalities give the same inclusion for the two edge intervals, since $m_\pm$ are fixed. The three starting-point intervals cover $[0,\Delta]$. Since $\Delta\ge1$, the middle estimate $p_0/2$ also gives a lower bound of $(p_0/2)\Delta^{-2}$. Taking the minimum of the three positive constants proves the claim.
    }{
        对中间区间, 当 $j\le n_\Delta$ 时有 $h_j\le c$, 故 $c\le S_j\le\Delta$ 蕴含 $h_j\le S_j\le\Delta+h_j$. 由 $h_{n_\Delta}\le c$ 及 $\frac25<\frac5{12}<\frac7{12}<\frac35$, 对充分大的 $\Delta$, 其中的终点区间包含于 $\mathcal A_\Delta(x)$ 所定义的区间. 对两侧区间, 上述终点不等式及 $m_\pm$ 固定也给出相同的包含关系. 三个起点区间覆盖 $[0,\Delta]$. 由于 $\Delta\ge1$, 中间区间的估计 $p_0/2$ 也给出下界 $(p_0/2)\Delta^{-2}$. 取三个正常数的最小值即得结论.
    }
\end{proof}
